\documentclass[11pt]{book}
\usepackage[paperwidth=145mm,paperheight=200mm,top=18mm,bottom=20mm,inner=20mm,outer=16mm]{geometry}
\usepackage{brumfiel}

% ---- local macros (hoist into book preamble) ----
% Unnumbered star footnote, as the book sets its five notes in this chapter.
% brumfiel.sty's \starnote emits a bare \footnotetext, which prints the
% footnote counter ("0"); this chapter's notes carry a bare "*" and no number.
\newcommand{\XVIstarnote}[1]{\begingroup\renewcommand{\thefootnote}{}%
  \footnotetext{*#1}\endgroup}
% In-text marker for those notes (a raised asterisk, as printed).
\newcommand{\XVIstar}{\textsuperscript{*}}
% ---- end local macros ----

% This chapter is unbroken prose in a narrow measure, with several long
% technical words; same allowance frontmatter.tex uses. Affects spacing only,
% never wording.
\emergencystretch=1.5em
\hyphenation{non-math-e-mat-i-cal non-Eu-clid-e-an epis-te-mol-o-gy
  math-e-ma-ti-cians char-ac-ter-is-tics grav-i-ta-tion-al phi-los-o-phers
  con-tra-dic-to-ry in-tel-lec-tu-al sub-sci-ences}

\begin{document}
\setcounter{chapter}{16}

\chapter*{\raggedright\Huge PHILOSOPHY OF\\MATHEMATICS}
\addcontentsline{toc}{chapter}{16.\ Philosophy of Mathematics}
\markboth{PHILOSOPHY OF MATHEMATICS}{}
\vspace{-1em}\noindent\rule{\linewidth}{0.6pt}\vspace{1em}

%========================================================= 16-1
\section*{16--1\quad WHAT IS PHILOSOPHY?}
\addcontentsline{toc}{section}{16--1\quad What is philosophy?}

The word ``philosophy'' is a combination of two Greek words and means
literally \emph{love of wisdom}. In the early Greek civilization, philosophy
was the study of all branches of knowledge. As man's learning increased
through the ages, certain disciplines grew until they split away from
philosophy and became separate areas for study. We no longer think of
medicine, law, mathematics, physics, chemistry, biology, economics, etc., as
parts of philosophy, although philosophy was the father of all these sciences.
That is, the first men who called attention to important concepts in
mathematics, physics, chemistry, economics, political science, etc., were
philosophers. It was philosophy that suggested how knowledge about these new
areas should be collected and organized.

Philosophy has always concerned itself with the unknown---with those
situations in which man is not sure of himself. Because of this, many of the
conclusions that philosophers reach cannot be checked by experiment, as we
check physical laws, and consequently many people have the mistaken belief
that time devoted to philosophical thought is largely wasted.

It is impossible to describe briefly and clearly just what the subject matter
of philosophy consists of today. Some of the branches into which philosophy is
subdivided will be listed here, but we make no pretense of completeness.

\emph{Moral philosophy}, or \emph{ethics}, investigates the nature of right
and wrong. It is concerned with distinguishing between good and bad conduct.
It raises questions like, ``Is an action to be judged good or bad according to
its effect or its intent?'' ``What are \emph{good} and \emph{evil}?'' ``How
shall we determine whether one action is better or worse than another?''

\emph{Religious philosophy} is the study of concepts common to religions. On
the other hand, \emph{theology} is concerned with the logical structure of a
single religion.

% NOTE (source oddity): the book jumps straight from "are constructed." into
% the two quoted questions, with no "It raises questions like," lead-in of the
% kind the moral-philosophy paragraph uses above. Verified at 300 dpi; it is
% printed that way. Do not "repair" this.
\emph{Political philosophy} does not consider details of government, but
rather the principles upon which various systems of government are
constructed. ``How much authority may the state exert over an individual?''
``Are there fundamental human rights, which no state can usurp? If so, what
are they?''

\emph{Metaphysics} (``beyond'' physics) is sometimes described as the study of
final or ultimate meanings. Much that is nonsense has been written in this
area, but it includes many deep and fascinating insights.

\emph{Epistemology} can be described briefly as the theory of knowledge. It
considers the general questions, ``How do we learn?'' and, ``How shall we
organize our knowledge?'' As a matter of fact, the philosophy of
\emph{mathematics} is just a branch of epistemology; so you may think of this
chapter as an introduction to a fragment of epistemology, which is itself a
fragment of philosophy.

It is impossible to discuss the philosophy of mathematics without considering
the philosophy of science, for mathematics is used as a tool in every science
to organize and interpret scientific knowledge. We shall mention how
mathematics thus often leads to the discovery of new scientific facts.

%========================================================= 16-2
\section*{16--2\quad THE TASK OF SCIENCE}
\addcontentsline{toc}{section}{16--2\quad The task of science}

All information about what we call the \emph{real} or \emph{external} world is
obtained from the five senses: sight, sound, taste, smell, and touch. Our
sense organs, for example the eye, send messages to the brain. A newborn baby
cannot do much with these messages, but as he grows older, he learns to
interpret them. We are always amazed at the accomplishments of a person who is
both blind and deaf (for example, Helen Keller), for it is difficult for us to
grasp how such a person can have much understanding of the world. What we are
inclined to forget is that all of us have learned how to interpret information
from our sense organs so as to make useful responses. As you read this page,
your eye receives impressions and sends messages to the brain. You have
learned to interpret and respond to these messages.

% NOTE (source oddity): "that is a touch" is printed without a comma after
% "is" (modern usage would set "that is, a touch"). Verified at 300 dpi.
% Reproduce as printed.
\emph{Every scientific observation is, at first, a sense impression, that is a
touch, smell, taste, sight, or sound, for some experimenter. It is the task of
scientists to understand, explain, and predict happenings in the world of our
sense impressions.}

Some philosophers have suggested that because our only knowledge of the
external world comes from sense perception, \emph{we cannot even prove that
the external world has any existence apart from our sense perception of it;
that perhaps the world that we believe we see, touch, and smell exists only as
things exist in a dream.}

% NOTE (source error, reproduce as printed): Berkeley is called a
% "17th-century" philosopher. He is 18th-century (1685--1753). The error is
% the book's --- see ch16-UNCERTAIN.md item 4. Do not silently correct.
The story is told that the great 17th-century philosopher Bishop Berkeley was
so convinced by this line of thought that one day, when he was walking and
meditating upon philosophical ideas, he encountered a large stone in his path
and so much doubted the actual existence of the stone that he drew back his
foot and gave it a vigorous kick. We have no record of the good bishop's
thoughts as he hobbled home.

%========================================================= 16-3
\section*{16--3\quad KINDS OF SCIENCE}
\addcontentsline{toc}{section}{16--3\quad Kinds of science}

The world is very complicated, and we have little understanding of it. To
simplify the task of understanding our environment, certain topics that seem
somewhat related are studied by themselves as separate sciences. For example,
biology is the study of living organisms, and biology is itself separated into
several subsciences, like zoology, the study of animal life, and botany, the
study of plant life. To understand all aspects of biology, it is necessary to
know some chemistry and physics. For this reason chemistry and physics are
regarded as more basic, or more fundamental, than biology. Of the sciences
concerned with the real world, the most fundamental is physics, which deals
with the properties of matter.

All sciences become \emph{quantitative} when they have progressed far enough.
This means that measurements are made and numbers are introduced into the
sciences. Those things which we measure, and to which we assign numbers, are
called \emph{quantities}. As soon as numbers are introduced into a science,
mathematics becomes necessary for its study. As a matter of fact, one needs
geometry in order to discuss the physics of motion of bodies, even before one
makes measurements. Even psychology and sociology, two of the newer sciences
that have only recently split away from philosophy, need statistics in order
to organize and interpret the huge number of facts with which they deal. Thus
all the sciences depend upon mathematics, so that it has been termed ``the
handmaiden of the sciences.''\XVIstar
% NOTE (source oddity): the footnote names Baltimore twice. Verified at
% 300 dpi --- the duplication is genuinely printed. Do not correct.
\XVIstarnote{E.~T. \textsc{Bell}, \emph{The Handmaiden of the Sciences}.
Baltimore: The Williams and Wilkins Co., Baltimore, 1937.}

Sometimes mathematics itself is classified as one of the sciences; and
sometimes one hears the phrase ``science \emph{and} mathematics.'' In any
event, call it what you may, mathematics is different from the other sciences
in that it is \emph{independent of experiments}. \emph{Mathematics is a
creation of the mind alone.}

%========================================================= 16-4
\section*{16--4\quad IS PHYSICS TRUE?}
\addcontentsline{toc}{section}{16--4\quad Is physics true?}

We cannot give an acceptable definition of physics in a few words, but it will
not take long to make the question of this section meaningful. You have
studied a little physics previously, so you should know that physics is
concerned with forces, motion, heat, sound, electricity, and much else. For
example, in that part of physics called mechanics, it is learned that if at
some instant we know the position and velocity of a ball, and if in addition
we know the forces acting upon it, then we can \emph{calculate} where it will
be at any later time.

You need not concern yourself here with the exact method used to make this
calculation. What is important is that \emph{given certain information, we can
predict what will happen}. Physics, and indeed every science, is just as
``true'' as its predictions are \emph{reliable}. As soon as a physical theory
fails to predict what will happen, then it becomes necessary to construct a
new theory that will predict more accurately. In physics, the final test of
``truth'' is whether the predictions we get from our theory actually agree
with our observations. As an example of this, Aristotle had written in his
\emph{Physics} that heavier bodies would fall more rapidly than lighter ones.
This prediction is a logical consequence of Aristotle's physical
\emph{theories}, or \emph{laws}. Every school child knows how Galileo
\emph{proved} these theories \emph{false} by dropping balls of various sizes
and weights from the leaning tower of Pisa. Every observer who was open-minded
enough to watch the experiment (and there were indeed some who were not)
became convinced that Aristotle's laws of motion were false and that it was
necessary to replace them by new laws that would come closer to predicting the
events that were actually observed.

The physical laws used by the physicist to make his predictions are actually
\emph{assumptions} and are expressed in terms of mathematical formulas. If the
mathematical calculations employed do not lead to correct predictions, this
does not mean that the mathematics is incorrect.\XVIstar\ Instead it means
that the physical laws do not agree well enough with reality, that the
\emph{physical assumptions} are at fault.
\XVIstarnote{Of course someone could make a computational mistake, but this is
the fault of the computer, not of the mathematics he is using.}

When the physicist defines physical concepts like force, weight, etc., and
then gives physical laws that tell how these concepts are related, he is
creating what we call a \emph{mathematical model}. What he is saying is that
\emph{if} there are things in the real world that correspond to these
concepts, and \emph{if} these things are related to one another according to
the rules called physical laws, \emph{then} mathematics tells us what will
happen in the real world.

If the predictions come out wrong, then the physicist must construct a better
mathematical model. Often it happens that the physicist cannot solve the
mathematical problems that arise within the mathematical model. Then he calls
upon the mathematician for help. In this way physics supplies mathematics with
interesting problems.

As an example of a simple mathematical model, consider our solar system. The
physicist thinks of the sun and planets as points that are tracing out curves
as they move through space. By using calculus and analytic geometry, these
curves are described by mathematical equations that predict how the planets
ought to move through the sky if the law of gravitation is valid.

% NOTE (source error, reproduce as printed): the book dates the
% Uranus/Neptune episode to "early in the 18th century". Adams and LeVerrier
% worked in the 1840s, and Uranus was not discovered until 1781. The error is
% the book's --- see ch16-UNCERTAIN.md item 4. Do not silently correct.
It is most interesting that early in the 18th century, when astronomers
calculated how the planets \emph{should} move according to the law of
gravitation, it was noticed that one of the planets, Uranus, did not follow
the exact path predicted for it by the mathematical model. The astronomers had
so much faith in the mathematical model that they became sure that there must
be an ``invisible'' planet, which was pulling Uranus off its course with its
gravitational attraction. The Englishman John C. Adams and the Frenchman
LeVerrier set about attempting to solve the problem of determining the path of
the ``invisible'' planet.\XVIstar\ Working with the mathematical model, they
computed where it \emph{ought} to be to affect the course of Uranus as it did.
When observers focused their telescopes upon that part of the heavens, they
saw the planet Neptune precisely where the mathematical model predicted that
it would be found. This is but one of the many dramatic illustrations of how
an \emph{abstract} mathematical model can lead us to the discovery of new
facts about our universe.
\XVIstarnote{A fascinating account of this work can be found in
J.~R.~Newman's \emph{The World of Mathematics}. Vol.~2, pp.~822--839.}

%========================================================= 16-5
\section*{16--5\quad IS MATHEMATICS TRUE?}
\addcontentsline{toc}{section}{16--5\quad Is mathematics true?}

If by this question we mean, ``Does the real world behave like the
mathematical models that we construct?'' then we are asking a nonmathematical
question. The \emph{truth} of mathematics is independent of the real world.
The \emph{usefulness} of one kind of mathematics or another does depend upon
the real world.

\emph{Pure mathematics is concerned with the logical consequences of postulate
systems.} In pure mathematics, postulate systems are studied primarily because
they are interesting, and often without regard to any future use. In the
history of mathematics, there are many postulate systems that have been
studied for their own sakes and only much later found to be applicable to the
external world. For example, the non-Euclidean geometry needed in the theory
of relativity was first studied many years before any physical application of
it was made.

\emph{Applied mathematics is concerned with those parts of pure mathematics
that are of greatest use in the science of today.} Scientific applications
often suggest the constructions of new postulate systems. Mathematics and
science thus serve as stimuli to each other.

We have not yet given a direct answer to the question, ``Is mathematics
true?'' What is ``truth,'' anyway? Each of us has some sort of feeling about
what ``truth'' means. Truth, for the physicist, is \emph{that which works},
and this attitude is characteristic of science. Truth, in practical affairs,
usually refers to \emph{whether some event actually took place}. Here, we
often take the word of others whom we trust. Truth, in spiritual matters,
depends upon \emph{faith}. In general, agreement among people regarding
spiritual beliefs is difficult to obtain, for there is no test by experience
that we can all agree upon.

When we say that a system of mathematical postulates is ``true,'' we mean
only, \emph{these postulates are free from contradiction}. That is, if it is
impossible to prove from our postulates that a statement is true, and then
also prove that the same statement is false, in this case and only in this
case do we call our postulates true. Remember that we do not assert in
mathematics that anything actually exists in the physical world. Rather, our
theorems have the form: \emph{If} a set of things has so-and-so properties,
\emph{then} it has such-and-such properties.

We have developed our geometry this year from statements that dealt with a set
of points that we called our plane. A mathematician is naturally pleased to
encounter in the physical world a set of objects that actually has the
properties he has postulated for one of his mathematical systems, for then he
can be sure that the set he finds will have other properties that he has
established by studying his abstract system. A natural question to ask at the
close of this year's study of geometry is whether there actually exists in the
real world a set of things that satisfies the postulates of our geometry. Does
the space we live in actually have the properties of the Euclidean space we
study in geometry?

The answer will probably not satisfy you. The answer is, \emph{we do not know.
Furthermore, no mathematician ever expects to know.} Our geometric plane
consists of \emph{infinitely} many points. \emph{It will never be possible to
determine whether our universe is finite or infinite in extent.} Although we
have no assurance that any part of the real world actually fits our geometric
postulates, experience has shown us that much of our geometric theory is
applicable to the world around us; and we can indeed use our geometry to make
very accurate predictions.

%========================================================= 16-6
\section*{16--6\quad FREEDOM FROM CONTRADICTION}
\addcontentsline{toc}{section}{16--6\quad Freedom from contradiction}

In this section we discuss a very deep question. When we have a postulate
system before us, how do we know that these postulates are
\emph{consistent}? How do we know that we shall never deduce from these
postulates, by correct logical reasoning, two contradictory conclusions? Let
us give an example, first given by Bertrand Russell,\XVIstar\ to illustrate
what we mean.
\XVIstarnote{An English philosopher who is interested in the foundations of
mathematics. At about the turn of the century, he wrote \emph{Principia
Mathematica} with A.~N. Whitehead. In the book they attempt to derive all
mathematics from logic alone.}

\emph{Russell's barber.} Let us assume that, in a certain small town, there is
but one barber. We assume that this barber shaves everyone in the town who
does not shave himself, and no one else.

We now ask who shaves the barber.

\emph{If} the barber does not shave himself, \emph{then} he is a person who
does not shave himself, and therefore he is shaved by the barber---himself.
Since this is impossible, we have proved that the statement \emph{``The barber
does not shave himself''} is false. Hence, the barber shaves himself. But, by
our assumptions, \emph{if} the barber shaves himself, \emph{then} he is one of
the men who is not shaved by the barber---again himself. This proves that the
statement \emph{``The barber shaves himself''} is false, and hence,
\emph{``The barber does not shave himself,''} is true. But now we have arrived
at a contradiction, having shown from our postulates that the statement
\emph{``The barber does not shave himself'' is both true and false.}

This example shows that in forming a set of postulates, we must be careful
that our postulates are consistent; that is, that they do not lead to
contradictory conclusions. When our postulate sets are very simple, it is
sometimes possible to prove that no contradiction can ever be reached by
logical reasoning. For more complicated postulate sets, this may not be
possible.

David Hilbert spent much of the latter portion of his life trying to prove the
consistency of certain interesting mathematical systems. At about the same
time that he wrote his \emph{Foundations of Geometry}, he showed that geometry
is consistent if arithmetic is consistent. He accomplished this proof by using
the ideas of Chapter 15. The details need not concern us here. All we need
know is that our geometry is consistent if arithmetic is. We should remark
that non-Euclidean geometries have also been proved to be consistent,
providing that arithmetic is. Of course this has nothing to do with the
question of whether the space we live in is Euclidean or non-Euclidean.

% NOTE: "0 = 1" is set inline; the exclamation mark after it is rhetorical,
% NOT a factorial. Do not typeset it as $0 = 1!$ in a way that reads as one.
Now, how do we know that arithmetic is consistent? The answer may again seem
disappointing, for the answer is, \emph{we do not know!} Perhaps we never will
know! There is no proof that we can never be led to an impossible conclusion
in arithmetic. For example, we do \emph{not} know that it is impossible to
\emph{prove}, by a valid chain of logical arguments, that $0 = 1$!

All mathematicians believe that arithmetic is consistent. This belief is based
upon faith, not proof. Of course, arithmetic is useful. It works, and helps us
reach conclusions that turn out to be correct. But here is the position in
which the scientist finds himself. If an inconsistency is discovered in
arithmetic, the postulates for arithmetic will have to be changed. The
possibility seems to be remote, so it does not worry working mathematicians
too much.

%========================================================= 16-7
\section*{16--7\quad MATHEMATICAL RESEARCH}
\addcontentsline{toc}{section}{16--7\quad Mathematical research}

Euclid's geometry is more than 2000 years old. The average person holds the
erroneous opinion that all of mathematics was discovered long ago, and that
today we study only old stuff. There are, however, many journals published
throughout the world that are devoted primarily to \emph{new} mathematics. In
fact, during the past fifty years, more new mathematics has been
invented\XVIstar\ than in all previous history! Mathematical progress is
stimulated by both intellectual curiosity and the needs of modern science.
Because of the mathematical nature of modern science and engineering,
mathematicians are in greater demand than ever before and command excellent
salaries.
\XVIstarnote{Is mathematics \emph{discovered} or \emph{invented}? The first
word suggests that mathematics is ``sitting around waiting to be found.'' The
second suggests that mathematics is created by the mind, that human activity
is essential.}

It is impossible to describe here the breadth and depth of current
mathematical research. As yet, we have not undertaken enough elementary
mathematics to appreciate such a description. However, in your study of plane
geometry, you have made a fine beginning by studying a single mathematical
system in some detail. The logical features of geometry are common to all
mathematics, so that you have before you one example of a postulate system and
some of its consequences.

%========================================================= 16-8
\section*{16--8\quad MATHEMATICS AND LIFE}
\addcontentsline{toc}{section}{16--8\quad Mathematics and life}

How may an understanding of the philosophy of mathematics affect you as an
individual in your everyday life? Mathematicians are much like any other
group. If you were to attend the International Congress of Mathematicians,
which convenes once every four years, you would see a group of people
differing widely in dress, speech, religious and political beliefs, and
physical appearance. But you would, if you were observant enough, see many
common characteristics. Everyone there would have a way to communicate with
everyone else---not through English, French, German, or Russian, but through
the universal symbolism of mathematics. You would soon realize that these
folk, while actively interested in the real world, were still far more
interested in the world of ideas.

You would see that each one present had a tremendous respect for the
personalities and intellectual capabilities of his neighbors.

The philosophy of mathematics speaks to those who will listen and assures them
that even more important than wresting food, shelter, and material riches from
the earth is the acquisition of knowledge and understanding.

%========================================================= Suggested Reading
\par\medskip
\begin{center}\bfseries SUGGESTED READING\end{center}
\par\nopagebreak\medskip

\noindent From \emph{The World of Mathematics}, by J.~R.~Newman.
\begingroup\raggedright
\begin{exlist}
\item \textsc{Philip} E.~B. \textsc{Jourdain}, ``The Nature of Mathematics.''
      Vol.~1, pp.~4--71.
\item J.~W.~N. \textsc{Sullivan}, ``Mathematics as an Art.''
      Vol.~3, pp.~2015--2021.
\item \textsc{Raymond} L. \textsc{Wilder}, ``The Axiomatic Method.''
      Vol.~3, pp.~1647--1667.
\end{exlist}
\endgroup

\end{document}
